\chapter{Testing and Verification}
\label{chap:testing}

Nothing in this project is user-facing, so testing means establishing that the
pipeline produces what it claims and that measured differences come from the
change under test. This chapter records how each part was verified and which
checks caught real errors.

\section{Verifying the reproduction}
\label{sec:verify-repro}

The first test was whether the unmodified implementation reproduces its
published result. Running it on EMNIST-10 for four hundred global rounds gave a
peak personalised accuracy of 96.45 per cent against the 96.49 reported in the
paper, a difference of 0.04 percentage points.

This matters more than a sanity check. The run carries the argument binding
defect, since it is present in the released code, and it still lands on the
published figure. That is evidence the defect was present when the paper's
results were produced, and it licenses treating later measurements as changes
to the method rather than symptoms of a broken build.

\section{Loader invariants}
\label{sec:loader-invariants}

Each loader is checked against six invariants before use.

No row appears in both the training and evaluation split for a client. Every
class present in the corpus appears in both the pooled training and pooled
evaluation sets. The pooled training features have approximately zero mean and
unit standard deviation, confirming the scaler was fitted and applied. No
non-finite values survive. The per-client, per-class split ratio lies within
0.70 and 0.80 of the intended 0.75. No client is empty and no team is empty.

The team check exists because it caught a real defect. One loader's sequential
assignment hardcodes its group boundaries to two specific client indices rather
than computing them, which at forty clients produces groups of eleven, twenty-six
and three and leaves the fourth team with no members. A team with no members
divides by zero when its sample proportion is computed.

\section{Errors the checks caught}
\label{sec:errors-caught}

Four errors reached a run before being caught, and the way each surfaced is
worth recording.

\subsection{A split that produced empty test classes}

The first temporal split ordered each client's rows by time and cut at
seventy-five per cent. Attacks in CICIDS2017 occupy contiguous windows, so this
gave the evaluation set whichever attack ran last and placed the others entirely
in training. Wednesday clients trained on 2,913 denial-of-service samples and
were evaluated on none.

The symptom was not an exception. It was that the personalised models scored
0.45 accuracy where a classifier that always predicts benign scores 0.82. A
model below the majority-class baseline indicates a setup fault rather than a
weak method, and that comparison is the check that caught it. The correction is
to split each class on its own timeline.

\subsection{A control flow error that passed the syntax check}

A guard inserted between a conditional branch and its trailing \texttt{else}
turned the construct into a \texttt{for}-\texttt{else}, which is valid Python.
The \texttt{else} then bound to the loop and raised on every call. Eight queued
runs each exited in three seconds. The parser had accepted the file.

The lesson recorded here is that a syntax check is not an execution check. Every
subsequent queue was preceded by a single short run of the same configuration.

\subsection{An undefined name from statement ordering}

A module-level assignment reading an environment variable was inserted above the
import that provides the module it calls. Forty-six queued runs failed
identically. Again the file parsed.

\subsection{A metric measured against a moving target}

Benign thinning was initially applied before the train and test split, so it
thinned both. The accuracy floor moved from 0.8171 to 0.7008, meaning
comparisons across thinning settings sat on different evaluation targets. The
correction, following the reference implementation, applies thinning to the
training split only, leaving the evaluation set at its natural class balance.
The floors are now constant across settings.

\section{Verifying the modifications}
\label{sec:verify-mods}

Each modification was verified by measuring the quantity it was supposed to
change.

For the argument binding fix, the optimiser's penalty weight was read directly
before and after: twenty before, matching the local iteration count, and the
supplied value after. The per-step pull changed from 0.2 to 0.05 as predicted.

For seeding, an identical configuration was run at two seeds and the results
compared. Before the change they were identical to four decimal places; after
it they differed.

For the synthetic split fix, the split was hashed across three separate
processes. Before the change the three hashes differed; after, they matched.

The clustering module was tested against synthetic data with known group
structure: twenty points drawn from five well-separated latent clusters. It
recovered them at an adjusted Rand index of 0.762 with equal group sizes, and
the trigger fired on separated data while holding off on near-identical data.

\section{Establishing what counts as a difference}
\label{sec:noise-floor}

A measurement is only interpretable against the variation of the measurement
itself. Two runs of an identical configuration, differing only in their
position in the experiment schedule, gave personalised macro F1 of 0.7318 and
0.7244, a difference of 0.0074.

Differences below roughly 0.01 are therefore not distinguishable from
run-to-run variation, and a threshold of 0.03 was adopted before results were
examined. Several comparisons fall below it and are reported as null rather than
as small effects. The largest such is 0.0016, which is an order of magnitude
inside the floor.

\section{What was not tested}
\label{sec:not-tested}

There is no unit test suite. The base implementation ships none, and adding one
would have meant characterising behaviour the project set out to question.
Verification here is at the level of measured invariants and reproduced
published figures, which suits an experimental study but would not suit
software intended for use.

The CUDA path is untested. All experiments ran on CPU and no GPU was available.
The device selection logic is exercised only in its CPU branch.
